% ============================================================
%  Peripheral Overview — STM32H7A3RITx
%  Generated from firmware.ioc + rtos.c
% ============================================================

\subsubsection{Peripheral Tables}
\label{sec:peripheral-tables}

The following tables document the complete hardware resource allocation of the
\mbox{STM32H7A3RITx} as configured in the firmware.
All timer frequencies are derived from the \SI{280}{\mega\hertz} CPU clock,
with APB2 running at \SI{140}{\mega\hertz} (timer clock \SI{280}{\mega\hertz} after the ×2 multiplier).

% ── Timers ────────────────────────────────────────────────
\paragraph{Timer Allocation}
\label{sec:timers}

\begin{table}[h]
	\centering
	\caption{Hardware timer allocation}
	\label{tab:timers}
	\small
	\begin{tabular}{llllll}
		\toprule
		\textbf{Timer} & \textbf{Channel} & \textbf{PSC / ARR} & \textbf{Frequency} & \textbf{Function} & \textbf{Notifies} \\
		\midrule
		TIM7   & —      & kernel default  & \SI{1}{\kilo\hertz}  & FreeRTOS timebase (SysTick alt.) & Kernel \\
		TIM13  & —      & 27999 / 499     & \SI{10}{\hertz}      & GPIO interface tick A (OPM)       & ControlIF, UserIF \\
		TIM14  & —      & 27999 / 199     & \SI{25}{\hertz}      & GPIO interface tick B (OPM)       & ControlIF, UserIF \\
		TIM15  & CH1    & — / 349         & \(\approx\)\SI{800}{\kilo\hertz} bit & WS2812 PWM bitstream & DMA1 Stream 4 \\
		TIM17  & —      & 50 / 24999      & \(\approx\)\SI{220}{\hertz}         & WS2812 animation tick             & UserInterfaceTask \\
		TIM1   & CH1N   & 50 / 49998      & \(\approx\)\SI{56}{\hertz}          & Fader LED 3 PWM (complementary)   & PWM hardware \\
		TIM12  & CH1/2  & 50 / 49998      & \(\approx\)\SI{56}{\hertz}          & Fader LED 1 \& 2 PWM              & PWM hardware \\
		\bottomrule
	\end{tabular}
\end{table}

\noindent
TIM13 and TIM14 operate in one-pulse mode (OPM), generating a single pulse per
trigger event rather than running continuously.
TIM15 is clocked such that its PWM duty cycle encodes WS2812 bit values;
the DMA transfer loads successive compare values into the capture/compare register
without CPU involvement.
TIM17's period-elapsed ISR issues a direct-to-task notification to
\texttt{UserInterfaceTask}, which calls \texttt{ws2812\_process} upon wakeup,
keeping all animation logic within the task context.

% ── DMA ───────────────────────────────────────────────────
\paragraph{DMA Channel Allocation}
\label{sec:dma}

\begin{table}[h]
	\centering
	\caption{DMA1 stream allocation}
	\label{tab:dma}
	\small
	\begin{tabular}{lllllll}
		\toprule
		\textbf{Stream} & \textbf{Request} & \textbf{Direction} & \textbf{Mode} & \textbf{Width} & \textbf{NVIC priority} & \textbf{Function} \\
		\midrule
		Stream 0 & SPI1\_RX & Periph $\to$ Mem & Circular & 16-bit & 5 & I2S audio RX (codec $\to$ MCU) \\
		Stream 1 & SPI1\_TX & Mem $\to$ Periph & Circular & 16-bit & 5 & I2S audio TX (MCU $\to$ codec) \\
		Stream 2 & ADC1     & Periph $\to$ Mem & Circular & 16-bit & 7 & Slide-potentiometer sampling \\
		Stream 3 & ADC2     & Periph $\to$ Mem & Circular & 16-bit & 8 & CV / V\slash Oct input sampling \\
		Stream 4 & TIM15\_CH1 & Mem $\to$ Periph & Normal  & 16-bit & 5 & WS2812 PWM compare values \\
		\bottomrule
	\end{tabular}
\end{table}

\noindent
All DMA streams operate without CPU involvement once initiated.
The I2S streams (0 and 1) use circular mode to implement a double-buffered audio
pipeline: the half-transfer and transfer-complete interrupts notify
\texttt{AudioTask} to process each half-block while the other half is being
transferred.
The WS2812 stream uses normal (non-circular) mode and is re-armed by the driver
for each LED frame.

% ── ADC ───────────────────────────────────────────────────
\paragraph{ADC Configuration}
\label{sec:adc}

Both ADC instances run in continuous conversion mode with DMA circular transfer,
so fresh conversion results are always available in memory without CPU polling.
Hardware oversampling (32$\times$ with a 5-bit right shift) raises the effective
resolution and reduces noise on all analogue inputs.

\begin{table}[H]
	\centering
	\caption{ADC channel assignment}
	\label{tab:adc}
	\small
	\begin{tabular}{lllll}
		\toprule
		\textbf{Instance} & \textbf{Channel} & \textbf{Pin} & \textbf{Label} & \textbf{Consumer} \\
		\midrule
		ADC1 & IN7  & PA7 & \texttt{FADER1\_IN} & UserInterfaceTask \\
		ADC1 & IN8  & PC5 & \texttt{FADER2\_IN} & UserInterfaceTask \\
		ADC1 & IN9  & PB0 & \texttt{FADER3\_IN} & UserInterfaceTask \\
		ADC1 & IN5  & PB1 & \texttt{FADER4\_IN} & UserInterfaceTask \\
		\midrule
		ADC2 & IN10 & PC0 & \texttt{V/OCT\_IN}  & ControlInterfaceTask \\
		ADC2 & IN11 & PC1 & \texttt{CV1\_IN}    & ControlInterfaceTask \\
		ADC2 & IN12 & PC2 & \texttt{CV2\_IN}    & ControlInterfaceTask \\
		ADC2 & IN13 & PC3 & \texttt{CV3\_IN}    & ControlInterfaceTask \\
		\bottomrule
	\end{tabular}
\end{table}

\noindent
ADC1 is configured as master and samples the four slide potentiometers.
ADC2 samples the four CV inputs including the V\slash Oct pitch input.
Both use an asynchronous clock of \SI{6}{\mega\hertz} derived from PLL3-R
(divided by 4 inside the ADC), with a sampling time of \num{64.5} cycles per channel.

% ── GPIO interrupts ───────────────────────────────────────
paragraph{GPIO and External Interrupts}
\label{sec:gpio}

\begin{table}[H]
	\centering
	\caption{GPIO interrupt (EXTI) assignment}
	\label{tab:exti}
	\small
	\begin{tabular}{lllll}
		\toprule
		\textbf{Pin} & \textbf{Label} & \textbf{Trigger} & \textbf{NVIC priority} & \textbf{Notifies} \\
		\midrule
		PA9  & \texttt{GATE1\_IN}   & Falling edge & 10 & UserInterfaceTask, ControlInterfaceTask \\
		PA8  & \texttt{GATE2\_IN}   & Falling edge & 10 & UserInterfaceTask, ControlInterfaceTask \\
		PC7  & \texttt{GATE4\_IN}   & Falling edge & 10 & UserInterfaceTask, ControlInterfaceTask \\
		PA10 & \texttt{BUTTON1\_IN} & Falling edge & 10 & UserInterfaceTask \\
		PA11 & \texttt{BUTTON2\_IN} & Falling edge & 10 & UserInterfaceTask \\
		\midrule
		PC9  & \texttt{GATE3\_IN}   & — (polled)   & —  & Polled by GPIO driver \\
		\bottomrule
	\end{tabular}
\end{table}

\noindent
GATE3 (PC9) is configured as a plain digital input without an EXTI line and is
therefore polled by the GPIO driver rather than interrupt-driven.
All other gate and button inputs generate falling-edge interrupts at NVIC
priority~10, which is below \texttt{configMAX\_SYSCALL\_INTERRUPT\_PRIORITY},
permitting the use of FreeRTOS API calls (specifically
\texttt{xTaskNotifyFromISR}) inside these handlers.

% ── Serial / audio buses ──────────────────────────────────
\paragraph{Serial and Audio Interfaces}
\label{sec:serial}

\begin{table}[H]
	\centering
	\caption{Serial and audio bus assignment}
	\label{tab:serial}
	\small
	\begin{tabular}{lp{3.5cm}p{3.5cm}p{5cm}}
		\toprule
		\textbf{Peripheral} & \textbf{Pins} & \textbf{Configuration} & \textbf{Function} \\
		\midrule
		I2S1 (SPI1) &
		PA4 WS, PA5 CK, PB4 SDI, PB5 SDO, PC4 MCK &
		Full-duplex master, \SI{48}{\kilo\hertz}, 16-bit &
		Bidirectional audio data to/from TLV320AIC3204 codec \\
		\addlinespace
		I2C1 &
		PB6 SCL, PB7 SDA &
		Timing register \texttt{0x20B0CCFF} &
		Codec register configuration and readback \\
		\addlinespace
		UART4 &
		PA0 TX, PA1 RX &
		Asynchronous &
		Debug output (auxiliary to SWO on PB3) \\
		\bottomrule
	\end{tabular}
\end{table}

\noindent
The I2S master clock (MCK) on PC4 is driven by PLL2-P at \SI{12.288}{\mega\hertz},
which is exactly $256 \times f_s$ at \SI{48}{\kilo\hertz} and allows the codec to
derive all internal clocks from a single external reference without its own PLL.

% ── Clock summary ─────────────────────────────────────────
\paragraph{Clock Summary}
\label{sec:clocks}

The STM32H7A3 uses three independent PLLs to derive the required clock domains.
PLL2 generates the I2S master clock at exactly $256 \times f_s = \SI{12.288}{\mega\hertz}$, allowing the codec to derive all internal clocks from this single reference without its own PLL.

The \SI{6}{\mega\hertz} ADC clock and 32$\times$ oversampling yield an effective per-channel rate of approximately \SI{610}{\hertz}, which is sufficient for slowly varying potentiometer inputs; however it does not allow for higher audio-rate CV tracking.
A future revision could increase the ADC clock and reduce oversampling on the CV input channels to approach sample rates suitable for pitch-accurate CV at audio rates.

\begin{table}[H]
	\centering
	\caption{Key derived clock frequencies}
	\label{tab:clocks}
	\small
	\begin{tabular}{lll}
		\toprule
		\textbf{Domain} & \textbf{Source} & \textbf{Frequency} \\
		\midrule
		CPU / AHB        & PLL1-P & \SI{280}{\mega\hertz} \\
		APB1 / APB2      & AHB / 2 & \SI{140}{\mega\hertz} \\
		Timer clock (APB×2) & — & \SI{280}{\mega\hertz} \\
		I2S / SPI123     & PLL2-P & \SI{12.288}{\mega\hertz} \\
		ADC              & PLL3-R / 4 & \SI{6}{\mega\hertz} \\
		HSE oscillator   & PH0/PH1 & \SI{24}{\mega\hertz} (external crystal) \\
		\bottomrule
	\end{tabular}
\end{table}

\subsubsection{DMA, and Memory Layout}
\paragraph{Memory Layout}

The STM32H7A3 provides several distinct memory regions with different performance  characteristics. 
The linker script assigns each region a specific role based on access speed and cache behaviour:

\begin{table}[H]
	\centering
	\begin{tabular}{llp{4.5cm}rr}
		\toprule
		\textbf{Region} & \textbf{Size} & \textbf{Usage} & \textbf{Used} & \textbf{\%} \\
		\midrule
		\texttt{FLASH} (Bank 1)      & \qty{1024}{kB} & Program code, constants (\texttt{.rodata}), load-time initialisers & \qty{100.6}{kB} & 4.8\% \\
		\texttt{FLASH} (Bank 2)      & \qty{1024}{kB} & Persistent settings and calibration data & --- & --- \\
		\texttt{DTCMRAM}             & \qty{128}{kB}  & \texttt{.data}, \texttt{.bss}, stack, heap, FreeRTOS stack & \qty{129.2}{kB} & 98.5\% \\
		\texttt{RAM}                 & \qty{1008}{kB} & Large static buffers (\texttt{.sram1}): audio sample buffers and reverb delay lines & \qty{960}{kB} & 93.0\% \\
		\texttt{RAM\_NONCACHEABLE}   & \qty{16}{kB}   & DMA buffers (\texttt{.dma\_buffer}) & \qty{1.1}{kB} & 6.8\% \\
		\texttt{ITCMRAM}             & \qty{64}{kB}   & Instruction TCM, reserved for latency-critical code & \qty{0}{kB} & 0\% \\
		\bottomrule
	\end{tabular}
	\caption{Memory region assignment and utilisation at build time.}
	\label{tab:memory_layout}
\end{table}

%\begin{listing}[H]
%	\begin{minted}{text}
	%	Memory region       Used Size    Region Size    %age Used
	%	DTCMRAM:            129160 B     128 KB         98.54%
	%	RAM_NONCACHEABLE:   1112 B       16 KB          6.79%
	%	RAM:                960000 B     1008 KB        93.01%
	%	ITCMRAM:            0 B          64 KB          0.00%
	%	FLASH:              100556 B     2 MB           4.79%
	%	\end{minted}
%
%	\caption{Memory region utilisation from GCC linker output.}
%	\label{lst:gcc-build-output}
%\end{listing}

Placing DMA buffers in non-cacheable memory is a gewollte design decicion: on the Cortex-M7, the D-cache and DMA transfers operate independently. 
If a DMA buffer resides in cached memory, the CPU may read stale data from the cache after a DMA write, or a 
DMA read may see outdated values that the CPU has not yet written back. 
ST's application note AN4839 explicitly recommends this approach: 
\textit{``Always better to use non-cacheable regions for DMA buffers''} \cite[pp.~6-7]{stmicroelectronicsLevel1Cache2018}. 
By marking the DMA region as non-cacheable in the MPU configuration, coherency is guaranteed without requiring explicit \texttt{SCB\_CleanDCache} / \texttt{SCB\_InvalidateDCache} calls around every transfer. \\

Buffer sizes for the tape buffers (Section~\ref{sec:tape_buffer}) and reverb delay lines (Section~\ref{sec:schroeder_memory_allocation}) are tuned to maximise SRAM utilisation: 
larger tape buffers extend the maximum recording time, while larger reverb delay lines allow for a more spacious room size. 
The GCC linker report (Table~\ref{tab:memory_layout}) confirms that both \texttt{RAM} (93\%) and \texttt{DTCMRAM} (98.5\%) are nearly fully utilised. 
Flash consumption is modest at under 5\%, leaving ample space for future firmware growth.

\newpage
\subsubsection{RTOS Task Structure}

The firmware employs FreeRTOS with three concurrent tasks organised by priority.
Figure~\ref{fig:freertos-flow-diagram} illustrates the task structure and signalling flow:

\begin{figure}[H]
    \centering
    \makebox[\textwidth][c]{%
        \includegraphics[width=1.15\textwidth]{../images/software_architecture/freertos_flow_diagram.drawio.pdf}
    }
    \caption{System Software Flow Diagram}
    \label{fig:freertos-flow-diagram}
\end{figure}

\paragraph{AudioTask}
The highest-priority \texttt{AudioTask} is responsible for all real-time audio processing.
After waiting on the binary semaphore \texttt{audioReadySemaphore}---which gates execution until the boot-time calibration routine has completed---it initialises the codec and audio engine, then enters a perpetual processing loop. Each iteration blocks on a direct-to-task notification issued by the DMA transfer-complete interrupt; upon wakeup, it fetches a consistent parameter snapshot via \texttt{param\_cache\_fetch}, processes one half-block of audio through the tape-player, exciter, and Schroeder reverb stages in sequence, writes the result back to the DMA output buffer, and drains any pending tape-transport commands from \texttt{tape\_cmd\_q} before yielding.

\paragraph{ControlInterfaceTask}
The mid-priority \texttt{ControlInterfaceTask} waits on task notifications from the ADC driver and, upon wakeup, calls \texttt{control\_interface\_process} to convert raw ADC readings into calibrated CV and pitch values.

\paragraph{UserInterfaceTask}
The lowest-priority \texttt{UserInterfaceTask} initialises the WS2812 LED driver and starts its animation timer. Then -- when both buttons are pressed during startup -- the boot-time calibration routine is executed, during which the audio and control tasks are explicitly suspended via \texttt{vTaskSuspend} to give the calibration sequence exclusive CPU access. 
Once calibration completes, \texttt{audioReadySemaphore} is released and the task thereafter responds to pot-movement and GPIO notifications to update the LED state and parameter mappings.

Animation updates are driven by a hardware timer ISR, which issues a 
direct-to-task notification on each tick; \texttt{UserInterfaceTask} calls 
\texttt{ws2812\_process} upon wakeup, keeping animation logic entirely within the task 
context and the ISR minimal.

\paragraph{Parameter Cache}
Parameter data propagated from the control and user-interface tasks to \texttt{AudioTask} bypasses FreeRTOS queuing entirely. 
\texttt{param\_cache} implements a lock-free shared-memory scheme in which lower-priority tasks write individual fields using atomic store operations, while \texttt{AudioTask} reads a consistent snapshot each cycle via
\texttt{param\_cache\_fetch}. 
This eliminates all blocking on the time-critical audio path and avoids priority inversion, following the established principle that real-time audio threads should never contend on a mutex~\cite{reinders2021}. 
Lock-free communication of this kind is well established in the real-time systems literature: Herlihy and Shavit provide a thorough theoretical grounding for wait-free and lock-free data structures~\cite{herlihy2012art}, while Butenhof gives a practical treatment of memory visibility and atomic operations in concurrent C programs~\cite{butenhof1997programming}.

\paragraph{Non-FreeRTOS Domain}
Outside the FreeRTOS task hierarchy, several interrupt-driven subsystems operate directly on the hardware. 
The I2S peripheral drives a circular DMA buffer that manages bidirectional audio data transfer between the codec and the MCU. 
A dedicated hardware timer controls the tick rate of the WS2812 animation system. 
GPIO interrupts notify \texttt{UserInterfaceTask} of button presses and gate-input events, and separately signal \texttt{ControlInterfaceTask} for gate-signal evaluation. 
The custom ADC driver raises ISRs on slide-potentiometer and \mbox{V/Oct} conversion-complete events and notifies the respective tasks when fresh data is available.

\subsubsection{Drivers and Peripherals}
% TODO: codec driver, ADC driver, CV input handling
\paragraph{Audio Codec Driver}

During hardware design, the versatile \textbf{TI TLV320AIC3204} stereo audio codec was chosen to handle analog-to-digital and digital-to-analog audio conversion.

The TLV320AIC3204 is a highly configurable audio codec controlled entirely through a register interface, exposing power management, input/output routing, gain stages, clocking, and pin multiplexing. It supports sample rates from 8~kHz to 192~kHz for both record and playback paths. The playback path provides programmable volume control, DSP filtering, and flexible DAC/analog mixing, while the record path supports both single-ended and differential input configurations with a digitally-controlled microphone preamplifier.

Clock generation is flexible: the internal clock can be derived from the \texttt{MCLK}, \texttt{BCLK}, or \texttt{GPIO} pins, or from an internal PLL whose input may likewise be sourced from any of these pins. The PLL accepts input clocks in the range of 512~kHz to 50~MHz, though its use is discouraged at the lowest power settings. The analog supply operates in the range of 1.5V--1.95V and the digital supply in the range of 1.26V--1.95V, with integrated LDOs accepting input voltages of 1.8V--3.6V. \cite{texasinstrumentsTlv320aic3204Datasheet}

A common MCLK of 12.288MHz is generated by STM32H7's PLL tree \ref{sec:clock}.
This is then used to derivate all needed sychronization signals via its internal clock dividers.

The TLV320AIC3204 is configured over I\textsuperscript{2}C using a thin register-access layer built on the STM32 HAL. The two primitive operations, \texttt{codec\_write\_reg} and \texttt{codec\_read\_reg}, each perform a blocking HAL transfer to the codec's 7-bit address (\texttt{0x18}), with any HAL error escalated immediately to the system \texttt{Error\_Handler}. Before initialisation, \texttt{codec\_i2c\_is\_ok} probes the bus with \texttt{HAL\_I2C\_IsDeviceReady} to confirm the device is accessible.

The \texttt{codec\_init} function executes a fixed register-write sequence to bring the codec into a known operating state. 
It begins with a software reset, then configures the clock dividers NDAC, MDAC, NADC, and MADC, whose values are selected at runtime from the \texttt{audio\_sample\_rate\_t} argument according to the relation $f_s = f_{MCLK} / (N \cdot M \cdot \text{OSR})$, with a 12.288~MHz master clock and a fixed OSR of 128. The I\textsuperscript{2}C interface is configured for 16-bit I\textsuperscript{2}S framing. 

The input impedance of the left and right MICPGA positive terminals is configured to \SI{40}{\kilo\ohm} by writing \texttt{0xC0} to Page~1, Register~\texttt{0x34} and Register~\texttt{0x34} respectively, selecting the \texttt{11} encoding of bits D7--D6 as defined in the register map \cite[p.~135]{TLV320AIC3204ApplicationReference2012}.

Subsequent writes on Page~1 enable the AVDD LDO, configure headphone soft-stepping, set the microphone PGA settling delay, and establish the reference charge time. 
The signal routing writes connect the DAC outputs to the line-output drivers LOL and LOR, and route the line inputs through the PGA to the ADC. 

The sequence concludes on Page~0 by powering up the DAC and ADC blocks, setting unity gain, and unmuting all signal paths.


\begin{listing}[h]
\begin{minted}{c}
/* sample-rate-dependent dividers: fs = MCLK / (N * M * OSR) */
switch (rate) {
	case SAMPLERATE_48KHZ: n_val = 1; m_val = 2;  break;
	case SAMPLERATE_24KHZ: n_val = 1; m_val = 4;  break;
	case SAMPLERATE_8KHZ:  n_val = 1; m_val = 12; break;
	default:               n_val = 1; m_val = 2;  break;
}
codec_write_reg(0x0b, 0x80 | n_val); /* NDAC - MSB powers up divider */
codec_write_reg(0x0c, 0x80 | m_val); /* MDAC */
codec_write_reg(0x0e, 0x80);         /* DAC OSR = 128 */
\end{minted}
\caption{TLV320AIC3204 clock divider configuration}
\label{lst:codec_clk}
\end{listing}

Full register details and their respective configuration values are documented in TI's \textit{TLV320AIC3204 Application Reference Guide} SLAA557, section~5.2 \cite[pp.~96--128]{TLV320AIC3204ApplicationReference2012}.

\paragraph{WS2812 RGB LED Driver}

The module uses four WS2812 addressable RGB LEDs for visual feedback, driven via a 
single-wire \qty{800}{kHz} PWM protocol where bit values are encoded as duty cycle 
widths ($\sim$\qty{0.35}{\micro s} for 0, $\sim$\qty{0.7}{\micro s} for 1). The 
bitstream is generated by a hardware timer in PWM mode with compare values precomputed 
into a double-buffered DMA buffer, allowing the next frame to be prepared while the 
current one transmits. On DMA completion, the output is held low for $>\qty{50}{\micro s}$ 
to satisfy the WS2812 reset requirement.

Each LED is managed by an independent state machine supporting four modes: off, static 
colour, trigger flash (used for gate indication), and animation. The animation engine 
runs at \qty{120}{Hz} driven by TIM17, linearly interpolating between keyframe stages 
at each tick. Looping and one-shot animations are supported, with a single queued 
\texttt{next\_anim} allowing smooth transitions from one-shot to idle animations.

Animation sequences are defined as statically allocated structs in 
\texttt{ws2812\_animations.c}, requiring no heap allocation. Defined animations include 
a white breathe (idle), red/blue mode indicators, a boot sequence, a chase pattern, and 
green/red flash sequences for confirmation and error feedback.

\subsubsection{Persistent Settings Storage}

User settings and calibration data are persisted to the second Flash bank 
(\texttt{FLASH\_BANK2}), leaving \texttt{FLASH\_BANK1} exclusively for program code. 
The \texttt{settings.c} module exposes a simple read/write API over two primitive 
operations: \texttt{flash\_write\_words} and \texttt{flash\_read\_words}.

Writing requires an erase-before-write cycle, as Flash sectors cannot be partially 
overwritten. The driver computes the affected sector range from the destination address, 
erases those sectors via \texttt{HAL\_FLASHEx\_Erase}, then programs the data in 
16-byte \textit{flashword} chunks — the minimum write granularity of the STM32H7 
Flash interface:

\begin{minted}{c}
	HAL_FLASH_Program(FLASH_TYPEPROGRAM_FLASHWORD, flash_addr + offset, block_ptr);
\end{minted}

The \texttt{SettingsData} struct is fixed at exactly \qty{64}{bytes} via explicit 
padding, and declared \texttt{\_\_attribute\_\_((aligned(16)))} to satisfy the 
STM32H7 flashword alignment requirement:

\begin{minted}{c}
	struct SettingsData {
		struct calibration_data calibration_data; // 36 bytes
		struct State state;                       //  4 bytes
		uint8_t padding[...];                     // 20 bytes
		uint32_t magic;                           //  4 bytes
	} __attribute__((aligned(16)));               // total: 64 bytes
\end{minted}

The fixed size ensures the struct always occupies exactly four flashwords, making the 
sector erase and write loop deterministic. The \texttt{magic} field at the end of the 
struct is written last and checked on startup — if it does not match 
\texttt{MAGIC\_NUMBER}, the application falls back to default calibration values, 
protecting against partially written or uninitialised Flash contents.

The calibration data itself stores the V/oct pitch scale and offset coefficients, 
per-channel CV offsets, and the potentiometer range endpoints — all floating-point 
values determined during the factory calibration procedure and required at boot before 
the audio engine starts.

Reading invalidates the DCache before accessing Flash 
(\texttt{SCB\_InvalidateDCache}), ensuring the CPU does not return a cached copy of 
previously read Flash contents after a write cycle. The Flash address is accessed via 
a \texttt{volatile} pointer to prevent the compiler from optimising away the reads.